\chapter*{Introduction}
\label{cap:introduction}
\addcontentsline{toc}{chapter}{Introduction}

\chapterquote{There is only one decisive victory: the last.}{Carl Von Clausewitz}

This Final Degree Project addresses the design and implementation of a digital collection
exploration engine that combines faceted filtering over metadata, transversal navigation
between related collections, and set operations over partial results. The proposal emerges in
response to a concrete problem in contemporary digital information consumption, motivated
below, and aims to give the user a tool that combines the precision of direct search, the
flexibility of tag-based filtering, and the expressive richness of hypermedia navigation.

This introductory chapter is organized in four blocks: the motivation that gave rise to the
project, the objectives pursued, the execution plan, and a description of how the rest of the
memoir is structured.

\section*{Motivation}

In recent times we have witnessed a major shift in the paradigm of digital consumption. The
quantitative accumulation of digital content has surpassed its qualitative threshold and now
demands a different way of engaging with it. The abundance of information generates what is
known as \textit{information overload}, in which any user requires an ever-growing amount of
time to navigate it. This becomes especially evident on digital platforms, particularly
multimedia streaming services (such as Netflix, Disney+, or HBO Max), where it is difficult to
find what to consume without investing a substantial amount of time in the process.

This shift in paradigm, and the consequent problem, is still being addressed with traditional
methods: direct search (by keywords), classical filtering by genre or media type, and
recommendation systems based on algorithms hidden from the consumer. These approaches are
limited and serve poorly in the new context of digital consumption. Direct search is
deterministic and allows no flexibility, especially as it assumes the user already knows
exactly what they are looking for. Traditional filtering is often incomplete and inconsistent,
relying on hidden and inscrutable taxonomies. Recommendation algorithms tend to serve more as
amplifiers of the \textit{filter bubble}, as they depend directly on consumption history or on
what the platform wishes to promote, ultimately stripping the user of their agency.

This problem, moreover, is not exclusive to user experience but also constitutes a
technological problem related to how information collections are designed and used. Relational
navigation models and \textit{collection navigation systems} bring a different approach in
which data is not represented as isolated elements but as entities connected through semantic
relationships. This approach makes it possible to traverse collections interactively, with more
expressive navigation mechanisms that let the user progressively explore a complex collection
by combining filters, relationships, and transitions between entities.

There is, therefore, a recognized need for navigation systems that accompany this paradigm
shift. This can only be achieved by providing users with richer and more flexible interactive
exploration tools. Such tools must combine the precision established by direct search with
greater expressiveness than classical filtering, without depriving the user of their agency.

\section*{Objectives}

This Final Degree Project starts from the original proposal of the supervising professor,
titled \textit{Application for the Processing of Digital Content Collections}, which envisioned
the design of a tool to transform metadata and content of objects in large digital collections,
with possibilities ranging from the recoding of descriptions to tagging assisted by large
language models. From that initial proposal we have narrowed the scope toward \textbf{interactive
exploration} of collections, preserving the central idea of a domain-agnostic system capable of
operating on heterogeneous collections, but concentrating the effort on the navigation
primitives that allow the user to traverse them expressively. This scoping decision is supported
by the experimental nature of the navigation engine we wanted to validate and by the desire to
deliver a functionally closed system at the end of the project; the integration of
language-model-assisted transformations over the same engine remains as a natural continuation,
captured later among the future work lines.

\subsection*{General objective}

To design and implement an interactive exploration application for digital collections that
combines faceted filtering over metadata, transversal navigation between related collections
through semantic references, and set operations over partial results, all under a
domain-agnostic data model that allows the system to be applied to collections of different
structural nature without substantial modifications.

\subsection*{Specific objectives}

From the general objective, we identify six specific objectives, formulated so that each can be
verified independently in the Testing and Validation chapter.

\begin{itemize}
\item \textbf{O1. Domain-agnostic data model.} Define a relational model that represents
collections, entities, descriptive metadata, and relationships among entities, without coupling
to the particularities of any specific domain.

\item \textbf{O2. Navigation engine with three primitives.} Design and implement the
exploration logic through a state machine that combines faceted filtering with include and
exclude modes, transversal navigation through \textit{links}, and set operations (specifically,
union) over partial results, maintaining history and coherent invariants between transitions.

\item \textbf{O3. Functional validation over distinct domains.} Empirically verify that the
model and the engine generalize to collections of structurally different nature, through the
integration of at least two public catalogs from distinct domains.

\item \textbf{O4. Exposure of the engine through a \textit{RESTful API}.} Design and develop
an HTTP services layer that exposes the engine's capabilities in a structured manner consumable
by external interfaces, with session management and error handling.

\item \textbf{O5. Interactive web interface.} Develop a client application that materializes
the engine's primitives in interface components understandable by users without prior
experience in faceted search, with special attention to readability, response immediacy, and
consistency of the exploration metaphor.

\item \textbf{O6. Empirical evaluation of usability.} Design and conduct a formal usability
study with external users following a standard protocol (\textit{think-aloud} accompanied by
SUS and SEQ questionnaires and direct observation), from which actionable evidence about the
quality of interaction and a prioritized catalog of incremental improvements are derived.
\end{itemize}

Deliberately out of scope, and reserved as continuation lines collected in the Conclusions
chapter, are: automatic metadata transformations assisted by language models, user persistence
and authentication, and distributed-scale deployment of the engine. These scope boundaries are
deliberate and allow the system to be closed coherently within the project's time horizon,
without compromising the central primitives that constitute the main contribution.

\section*{Work plan}

We adopted an agile methodology, based for the most part on the \textbf{Scrum} framework,
adapted to a paired project. The choice of this method stems from the exploratory and
innovative nature of the project, since the requirements and the navigation engine required
iterative refinement and optimization.

Development was organized through \textit{sprints} (work cycles of two to three weeks),
defining an initial \textit{Product Backlog} with the fundamental system requirements. Periodic
follow-up meetings (\textit{Sprint Reviews}) allowed us to evaluate each software increment and
adapt planning according to the problems and obstacles encountered, jointly with supervision
from the professor.

Regarding the macro-planning, the project was structured into two sequential development
phases. The first phase, devoted to investigation and the Java prototype (September 2025 --
January 2026), aimed at validating the technical viability of the navigation engine and the
algorithms involved by deliberately postponing the implementation of a graphical interface and
centralizing the effort on the technical efficacy, efficiency, and robustness of the system's
components. The second phase, dedicated to web development and integration (January 2026 --
March 2026), shifted the effort toward exposing the engine's functionality through a
client-server architecture and developing the user interface, including the design of the
\textit{RESTful} controller layer using Java Servlets, the implementation of the web frontend,
and the integration testing with the full datasets. The detailed activity breakdown, the
chronogram, and the Gantt diagram are reproduced in the Spanish version of this chapter and
are not duplicated here.

\section*{Organization of the memoir}

The remainder of the memoir is structured into five numbered chapters, a section on
\textit{Personal Contributions}, and three appendices, organized following the classical
progression of framing, design, implementation, validation, and closing.

The \textbf{State of the Art} chapter situates the work within the theoretical framework of
tag-based navigation and structured hypermedia models, and reviews the related works the
system most directly engages with, particularly the Clavy platform and recent proposals based
on large language models. It identifies which project decisions are inherited from prior
contributions and which deliberately take a different path.

The \textbf{Architecture and Methodology} chapter presents the organization of the software
engineering process followed, with the Scrum adaptation for a small team, and the global
system architecture in terms of layers, modules, and contracts between them. It closes with a
description of the persistence layer and the data access strategy over PostgreSQL.

The \textbf{Development and Implementation} chapter dives into the concrete components that
materialize the engine's primitives. It goes through, in this order, the relational data
model, the representation of the exploration context (the state machine), classical and
relational filtering, transversal navigation through \textit{links}, set operations over
partial results, the web client development, and the obstacles encountered together with the
solutions applied.

The \textbf{Testing and Validation} chapter gathers the system validation from two
complementary perspectives. The first one, functional, over the two selected test domains
(TMDB and DBLP). The second one, which constitutes the bulk of the chapter, is a formal
usability evaluation with external users following a \textit{think-aloud} protocol accompanied
by SUS and SEQ questionnaires. Quantitative results, qualitative findings, derived interface
iterations, and validated use cases are reported, and the chapter closes with an explicit
enumeration of the study's limitations.

The \textbf{Conclusions and Future Work} chapter synthesizes the degree of fulfilment of the
objectives stated in this chapter, acknowledges the limits of the delivered system, and
enumerates the five natural continuation lines identified from the state of the work. The
\textit{Personal Contributions} section that follows documents the distribution of work between
the two authors.

Three appendices complete the memoir. The first reproduces the complete usability evaluation
script, so the study can be replicated unchanged by third parties. The second describes the
integration and ingestion of the external data sources (TMDB and DBLP), with detailed coverage
of the capture and transformation \textit{scripts}. The third documents the system deployment
over containers and gathers the operational decisions that allow the execution environment to
be reproduced.
